% ===== PRÉAMBULE CANONIQUE — à coller VERBATIM en tête de CHAQUE .tex =====
\documentclass[11pt,a4paper]{article}
\usepackage[utf8]{inputenc}\usepackage[T1]{fontenc}\usepackage{lmodern}
\usepackage[french]{babel}
\usepackage{amsmath,amssymb,mathtools}\usepackage{icomma}
\usepackage[version=4]{mhchem}          % \ce{...} : équations & formules
\usepackage{siunitx}\sisetup{locale=FR} % \qty{}{}, \si{}, \ang{}
\usepackage{chemfig}                    % \chemfig{...} : structures développées
\usepackage{xcolor}\usepackage{colortbl}
\usepackage{tikz}\usepackage{pgfplots}\pgfplotsset{compat=1.18}
\usepackage[most]{tcolorbox}\usepackage{enumitem}\usepackage{array,booktabs}
\usepackage[a4paper,margin=2.2cm]{geometry}
\usetikzlibrary{arrows.meta,calc,angles,quotes,positioning,patterns,backgrounds,shapes.geometric}
\definecolor{primary}{RGB}{222,49,99}\definecolor{primarydark}{RGB}{150,22,55}
\definecolor{defcol}{RGB}{0,114,178}\definecolor{methcol}{RGB}{52,92,148}
\definecolor{excol}{RGB}{0,135,90}\definecolor{attncol}{RGB}{200,40,55}
\tcbset{boxbase/.style={enhanced,breakable,boxrule=0.9pt,arc=2.5pt,
  left=3mm,right=3mm,top=2.3mm,bottom=2.3mm,fonttitle=\bfseries,coltitle=white}}
% --- boîtes TUTORIEL (noms EXACTS, ne pas renommer) ---
\newtcolorbox{cle}[1][]{boxbase,colback=defcol!6,colframe=defcol,
  title={\textsc{À retenir}\ifx\relax#1\relax\relax\else\textmd{ : \itshape #1}\fi}}
\newtcolorbox{methode}[1][]{boxbase,colback=methcol!7,colframe=methcol,sharp corners,
  title={\textsc{Méthode}\ifx\relax#1\relax\relax\else\textmd{ --- \itshape #1}\fi}}
\newtcolorbox{exemplebox}[1][]{boxbase,colback=excol!5,colframe=excol,
  title={\textsc{Exemple}\ifx\relax#1\relax\relax\else\textmd{ : \itshape #1}\fi}}
\newtcolorbox{piege}[1][]{boxbase,colback=attncol!6,colframe=attncol,
  title={\textsc{Piège}\ifx\relax#1\relax\relax\else\textmd{ : \itshape #1}\fi}}
\newtcolorbox{remarque}[1][]{boxbase,colback=black!4,colframe=black!45,
  title={\textsc{Remarque}\ifx\relax#1\relax\relax\else\textmd{ : \itshape #1}\fi}}
% --- boîtes EXERCICE / CORRIGÉ (noms EXACTS, auto-comptées) ---
\newtcolorbox[auto counter]{exo}[1][]{boxbase,colback=primary!4,colframe=primary,
  title={\textsc{Exercice}~\thetcbcounter\ifx\relax#1\relax\relax\else\textmd{ --- \itshape #1}\fi}}
\newtcolorbox[auto counter]{corrige}[1][]{boxbase,colback=excol!4,colframe=excol,
  title={\textsc{Corrigé}~\thetcbcounter\ifx\relax#1\relax\relax\else\textmd{ --- \itshape #1}\fi}}
\newcommand{\R}{\mathbb{R}}\newcommand{\N}{\mathbb{N}}\newcommand{\Z}{\mathbb{Z}}\newcommand{\ds}{\displaystyle}
% =========================================================================
\begin{document}
\sisetup{per-mode=symbol}
\begin{exo}[équations de dissolution avec ions polyatomiques]
Écris l'équation de dissolution de chacun de ces sels. Attention aux parenthèses et aux coefficients.
\begin{enumerate}[label=\alph*)]
  \item \ce{Ag2SO4}
  \item \ce{Ca3(PO4)2}
  \item \ce{Mg(OH)2}
  \item \ce{(NH4)2CO3}
  \item \ce{Al2(SO4)3}
\end{enumerate}
\end{exo}

\begin{exo}[calculer $M$ puis convertir $s_m\to s$]
Calcule d'abord la masse molaire $M$ du sel, puis sa solubilité molaire $s=s_m/M$.
\emph{Données :} $M(\ce{Pb})=\qty{207.2}{\gram\per\mole}$, $M(\ce{Br})=\qty{79.9}{\gram\per\mole}$,
$M(\ce{Ag})=\qty{107.9}{\gram\per\mole}$, $M(\ce{Cr})=\qty{52.0}{\gram\per\mole}$,
$M(\ce{O})=\qty{16.0}{\gram\per\mole}$.
\begin{enumerate}[label=\alph*)]
  \item \ce{PbBr2} : $s_m=\qty{4.84}{\gram\per\litre}$.
  \item \ce{Ag2CrO4} : $s_m=\qty{2.16e-2}{\gram\per\litre}$.
\end{enumerate}
\end{exo}

\begin{exo}[quel volume d'eau pour tout dissoudre ?]
La solubilité $s$ étant le nombre maximal de moles dissoutes par litre, détermine le volume $V$ d'eau
pure nécessaire pour dissoudre la quantité $n$ indiquée.
\begin{enumerate}[label=\alph*)]
  \item \ce{CaC2O4} : $s=\qty{3.0e-6}{\mole\per\litre}$, $n=\qty{6.0e-3}{\mole}$.
  \item \ce{CaCO3} : $s=\qty{6.9e-5}{\mole\per\litre}$, $n=\qty{6.9e-3}{\mole}$.
\end{enumerate}
\end{exo}

\begin{exo}[masse dissoute dans un volume donné]
À saturation, quelle masse $m$ de sel est dissoute dans le volume $V$ d'eau indiqué ?
On rappelle $m = s\times V\times M$.
\begin{enumerate}[label=\alph*)]
  \item \ce{BaSO4} : $s=\qty{1.0e-5}{\mole\per\litre}$, $M=\qty{233.4}{\gram\per\mole}$, $V=\qty{2.0}{\litre}$.
  \item \ce{CaF2} : $s=\qty{2.1e-4}{\mole\per\litre}$, $M=\qty{78.1}{\gram\per\mole}$, $V=\qty{0.25}{\litre}$.
\end{enumerate}
\end{exo}

\begin{exo}[concentrations ioniques à partir de $s$]
Connaissant la solubilité molaire $s$, calcule les concentrations des deux ions à saturation.
\begin{enumerate}[label=\alph*)]
  \item \ce{Ag2SO4} (type $2{:}1$) : $s=\qty{1.5e-2}{\mole\per\litre}$.
  \item \ce{Ca3(PO4)2} (type $3{:}2$) : $s=\qty{1.0e-7}{\mole\per\litre}$.
  \item \ce{Fe(OH)3} (type $1{:}3$) : $s=\qty{2.0e-10}{\mole\per\litre}$.
\end{enumerate}
\end{exo}
\end{document}
